\documentclass[10pt,a4paper]{article}
\usepackage[margin=2cm]{geometry}
\usepackage{fancyvrb,enumitem}
\fvset{fontsize=\scriptsize,frame=single}
\title{Lab Report -- Goal-Based Warehouse Agent with an LLM}
\author{Claude (AI assistant, Anthropic)}\date{}
\begin{document}\maketitle

\section*{Task 1: Understanding the problem}
\begin{enumerate}[leftmargin=*]
\item \textbf{Environment:} a $7\times21$ grid warehouse; \texttt{\#} cells are impassable shelving, \texttt{.} cells are free; fully observable, static, deterministic, discrete, single agent.
\item \textbf{Goal:} move the vehicle from \texttt{S} (row 1, col 1) to \texttt{G} (row 1, col 19) along a collision-free path.
\item \textbf{Actions:} Up, Down, Left, Right; each moves one cell and is legal only if the target cell is inside the grid and not \texttt{\#}.
\item \textbf{Information maintained:} the map (model of the world), the current position (state), the goal position, and the search bookkeeping (frontier, visited set, parent links) needed to build the plan.
\item \textbf{Goal-based, not simple reflex:} a reflex agent maps the current percept to an action by fixed rules (e.g.\ ``if free on the right, go right'') and would get stuck in dead ends; this agent holds an explicit goal and considers the future consequences of action sequences (search) before acting.
\end{enumerate}
\textit{Think about it (twice as large):} BFS is still correct and optimal, but its frontier grows with the number of cells (memory $O(|cells|)$ and time roughly 4$\times$ for a doubled area), so larger or weighted maps call for A* with a heuristic (Manhattan distance), which expands far fewer nodes. Further difficulties: moving obstacles or other vehicles, partial observability, non-uniform move costs, and replanning when the map changes.

\section*{Task 2: Agent design}
\begin{itemize}[leftmargin=*]
\item \textbf{Environment:} warehouse grid. \textbf{State:} current cell $(r,c)$. \textbf{Goal:} cell \texttt{G}. \textbf{Actions:} Up/Down/Left/Right. \textbf{Decision-making component:} BFS planner that searches the transition model for an action sequence reaching the goal.
\end{itemize}
\begin{Verbatim}
        +-------------------------- AGENT --------------------------+
        |                                                           |
 percept|  +-----------+   +---------------------+   +-----------+  |
 (map,  |  | State /   |   | Model: how actions  |   | Goal:     |  |
 position) | position  |-->| change position;    |<--| reach G   |  |
 ------>|  | (r, c)    |   | walls block moves   |   +-----------+  |
        |  +-----------+   +----------+----------+                  |
        |                             | "what happens if I do A?"   |
        |                  +----------v----------+                  |
        |                  | Decision maker:     |                  |
        |                  | BFS search -> plan  |                  |
        |                  +----------+----------+                  |
        +-----------------------------|-----------------------------+
                                      | actions (Up/Down/Left/Right)
                                      v
                              WAREHOUSE (environment)
\end{Verbatim}

\section*{Task 3: Prompt engineering}
\textbf{Prompt used} (the lab's suggested prompt, plus the exact map and testing requirements):
\begin{Verbatim}
Write a well-documented Python program implementing a goal-based agent for the
warehouse navigation problem shown below.
Map (S start, G goal, # obstacle, . free; moves Up/Down/Left/Right, one cell each):
#####################
#S....#............G#
#.##....##########..#
#....##.............#
#.######.###.#.###..#
#........#..........#
#####################
The program should
- represent the warehouse as a two-dimensional grid;
- determine a collision-free path from S to G;
- avoid all obstacles;
- print either the path found or a suitable message if no path exists;
- explain the search algorithm that has been chosen and why it is appropriate.
Also write a small test file that checks the path is valid, that its length is
shortest, and that an unreachable goal prints the no-path message.
\end{Verbatim}
I (Claude) wrote the code from this prompt.

\textbf{Test results.} \texttt{warehouse\_agent.py} ran without errors; it found a 20-move path (\texttt{Right x3, Down, Right x3, Up, Right x12}) which avoids all \texttt{\#}. \texttt{test\_agent.py} confirmed: path ends on \texttt{G}, never enters an obstacle, every step is one cell; its length equals an independently computed shortest distance (20); and a blocked map prints ``No collision-free path exists''. Path drawn on the map:
\begin{Verbatim}
#####################
#S***.#************G#
#.##****##########..#
#....##.............#
#.######.###.#.###..#
#........#..........#
#####################
\end{Verbatim}

\textbf{Questions.}
\begin{enumerate}[leftmargin=*]
\item \textbf{Working on first attempt?} Yes: the first generated version ran and passed all tests, so no fix iteration was needed. It was not edited; it was verified with the tests above rather than by trusting the printed path.
\item \textbf{Improving the prompt anyway:} it helped to give the exact map, name the required tests (validity, optimality, no-path case), and, if needed, specify the output format, the language version, and ``do not use external libraries''. A vaguer prompt (``write a program'') leaves the algorithm, map format and output undefined.
\item \textbf{Algorithm:} Breadth-First Search over grid cells, with a visited/parent dictionary and path reconstruction.
\item \textbf{Why:} all moves cost 1 and the state space is tiny, so BFS is complete and returns a shortest path; it is simple and needs no heuristic. (A* would also work and scales better; Dijkstra would matter only with varying costs.) The prompt's words ``determine a collision-free path'' on an unweighted grid point to BFS.
\end{enumerate}

\section*{Strengths and limitations of LLM-assisted development}
Strength: fast, well-structured, documented code from a precise specification. Limitation: the printed answer looks plausible whether or not it is correct, so independent tests (obstacle check, optimality check, unreachable case) are what establish correctness; the LLM also picks the algorithm implicitly, so the human must check the choice suits the problem (and that it would scale).

\section*{Code}
\subsection*{warehouse\_agent.py}\VerbatimInput{warehouse_agent.py}
\subsection*{test\_agent.py}\VerbatimInput{test_agent.py}
\end{document}
